%% Chapter 3 -----------------------------------------------------------------
\chapter{The Main Window}
\label{ch:main-window}
\index{main window}

The main window is where dynamic spectra are displayed, processed and measured, and from
where every other tool is opened. This chapter describes its layout, toolbar, menus and
status bar. The sidebar panels are described in Chapter~\ref{ch:sidebar}.

\section{Layout}
\label{sec:layout}

The main window (Figure~\ref{fig:main-window}) consists of the regions listed in
Table~\ref{tab:main-regions}.

\screenshot{main-window/main_window}{Main window of the \appname{} before a file is loaded.}{fig:main-window}{Main window with the regions labeled.}

\begin{table}[!htbp]
  \centering
  \caption{Regions of the main window.}
  \label{tab:main-regions}
  \small
  \begin{tabularx}{\linewidth}{@{}P{0.24\linewidth}L@{}}
    \toprule
    \thead{Region} & \thead{Description} \\
    \midrule
    Menu bar & Access to every function, including those not on the toolbar (Section~\ref{sec:menus}). On macOS the menu bar is at the top of the screen. \\
    Toolbar & Icon buttons for the most frequent actions (Section~\ref{sec:toolbar}). \\
    \gui{Cite this Software}\index{Cite this Software} & Opens the citation dialog (Section~\ref{sec:cite-dialog}). \\
    Sidebar & Collapsible panels for the timeline, background subtraction, display thresholds, units, axis scale, graph properties, the analysis summary and the ruler readout (Chapter~\ref{ch:sidebar}). \\
    Main plotting area & The dynamic spectrum with its color bar. Interactive tools act here. \\
    Status bar & Messages on the left; the update status and the live cursor readout on the right (Section~\ref{sec:status-bar}). \\
    \bottomrule
  \end{tabularx}
\end{table}

A narrow arrow handle on the left edge of the plotting area hides or shows the sidebar to
give the plot more room. The divider between the sidebar and the plot can also be dragged.

\section{Toolbar}
\label{sec:toolbar}
\index{toolbar}

The toolbar (Figure~\ref{fig:toolbar}) contains 16 tools, described in
Table~\ref{tab:toolbar}. Hovering over a tool shows its name. Tools that need loaded or
processed data remain disabled until they can be used.

\screenshot{main-window/toolbar}{The main toolbar with its tools numbered.}{fig:toolbar}{Toolbar with numbered callouts 1--16.}

\begin{xltabular}{\linewidth}{@{}c P{0.28\linewidth} L P{0.14\linewidth}@{}}
  \caption{Toolbar tools.}\label{tab:toolbar}\\
  \toprule
  \thead{No.} & \thead{Tool} & \thead{Function} & \thead{Shortcut} \\
  \midrule
  \endfirsthead
  \toprule
  \thead{No.} & \thead{Tool} & \thead{Function} & \thead{Shortcut} \\
  \midrule
  \endhead
  \callout{1} & Open / Load & Load one or more CALLISTO FITS files from disk (Section~\ref{sec:opening-files}). & \kbd{Ctrl+O} \\
  \callout{2} & Download & Open the e-CALLISTO FITS Downloader (Chapter~\ref{ch:downloader}). & \\
  \callout{3} & Export & Export the displayed spectrum as a figure (Section~\ref{sec:export-figure}). & \kbd{Ctrl+E} \\
  \callout{4} & Export as FITS & Export the processed spectrum as a FITS file (Section~\ref{sec:export-fits}). & \kbd{Ctrl+F} \\
  \callout{5} & Save Project & Save the complete workspace as an \file{.efaproj}\index{efaproj@.efaproj file} project (Section~\ref{sec:projects}). & \kbd{Ctrl+S} \\
  \callout{6} & Undo & Undo the last processing change. & \kbd{Ctrl+Z} \\
  \callout{7} & Redo & Redo the last undone\index{undo} change. & \kbd{Ctrl+Shift+Z} \\
  \callout{8} & Estimate Drift Rate & Estimate the drift rate from points clicked along a burst (Section~\ref{sec:drift}). & \\
  \callout{9} & Isolate Burst & Draw a freehand mask around a burst (Section~\ref{sec:isolation}). & \\
  \callout{10} & Plot Maximum Intensities & Open the Maximum Intensities window (Section~\ref{sec:max-intensities}). & \\
  \callout{11} & Rectangular Zooming & Drag one rectangle to zoom into the plot. Available while navigation is locked (Section~\ref{sec:navigation}). & \\
  \callout{12} & Lock zooming and panning & Disable scroll-wheel zooming and panning, freezing the view and enabling rectangle zoom. & \\
  \callout{13} & Unlock zooming and panning & Re-enable scroll-wheel zooming and panning. & \\
  \callout{14} & Reset Selection & Clear the isolated-burst selection and drift markers. & \\
  \callout{15} & Reset to Raw & Discard all processing and show the raw spectrum. & \\
  \callout{16} & Reset All & Clear the workspace: data, processing, selections and analysis overlays. & \\
  \bottomrule
\end{xltabular}

\begin{caution}
  \gui{Reset All}\index{Reset All} clears the whole workspace. Save a project first if you want to keep
  the current state.
\end{caution}

\section{Plotting area and renderers}
\label{sec:renderers}
\index{renderer}\index{hardware acceleration}

The plotting area shows the dynamic spectrum with frequency on the vertical axis, time on
the horizontal axis and intensity as color. The plot title shows the source file name
and the processing state, for example \gui{-Background Subtracted} or
\gui{-Isolated Burst}.

Two renderers are available and are selected with \menu{View > Mode}\index{renderer!Classic and Modern}:
\begin{description}
  \item[Modern]\index{renderer!Modern} uses a hardware-accelerated canvas (PyQtGraph\index{PyQtGraph}) for fast zooming and
    panning of large or combined spectra. This is the default.
  \item[Classic] uses Matplotlib\index{Matplotlib}.
\end{description}
\menu{Processing > Hardware Acceleration > Enable} switches the accelerated live canvas on
or off. Both renderers support the same tools, annotations and overlays, and exported
figures look the same whichever renderer is in use (Section~\ref{sec:export-figure}).

\section{Menu bar}
\label{sec:menus}
\index{menus}

Table~\ref{tab:menus} summarizes the menus of the main window. The tables that follow list
every command, with its shortcut and the section where it is described.

\begin{table}[!htbp]
  \centering
  \caption{Menus of the main window.}
  \label{tab:menus}
  \small
  \begin{tabularx}{\linewidth}{@{}P{0.18\linewidth}L@{}}
    \toprule
    \thead{Menu} & \thead{Main functions} \\
    \midrule
    \menu{File} & Open and combine FITS files; recent files and projects; open, save and recover projects; project reports; figure, FITS, provenance and analysis-log export. \\
    \menu{Edit} & Undo and redo; extend the timeline with adjacent observations; reset to raw data; reset the workspace. \\
    \menu{Download} & Launch the e-CALLISTO FITS Downloader. \\
    \menu{Solar Events} & CME catalog, GOES X-ray and proton flux, geomagnetic indices, SunPy Explorer, radio-burst downloaders, time-window synchronization, GOES overlay and the SWAVES panel. \\
    \menu{View} & FITS header, display range and its presets, view configurations, Multi-Station Comparison, theme and rendering mode. \\
    \menu{Analysis} & Solar Image Analysis, GCS CME Fitting, maximum intensities, Type~II band-splitting, light curves and the ruler. \\
    \menu{Processing} & Hardware acceleration, RFI cleaning, annotations, presets, automatic outlier cleaning, restored analysis sessions and batch processing. \\
    \menu{About} & Update check, bug report and version information. \\
    \menu{Help} & The built-in user guide. \\
    \bottomrule
  \end{tabularx}
\end{table}

\begin{xltabular}{\linewidth}{@{}P{0.36\linewidth} L P{0.17\linewidth}@{}}
  \caption{File menu.}\label{tab:menu-file}\\
  \toprule
  \thead{Command} & \thead{Function} & \thead{Shortcut, see} \\
  \midrule\endfirsthead
  \toprule \thead{Command} & \thead{Function} & \thead{Shortcut, see} \\ \midrule\endhead
  \menu{Open} & Load one file, or combine a compatible multi-file selection. & \kbd{Ctrl+O}, §\ref{sec:opening-files} \\
  \menu{Recent Files} & Reopen one of the last ten FITS selections. & §\ref{sec:recent-files} \\
  \menu{Combine FITS Files...} & Validate, preview and import time, frequency or time~\(\times\)~frequency combinations. & §\ref{sec:combine-dialog} \\
  \menu{Open Project...} & Open a saved \file{.efaproj} project. & \kbd{Ctrl+Shift+O}, §\ref{sec:projects} \\
  \menu{Recent Projects} & Reopen one of the last ten projects. & §\ref{sec:projects} \\
  \menu{Save Project} & Save the workspace to the current project file. & \kbd{Ctrl+S} \\
  \menu{Save Project As...} & Save the workspace to a new project file. & \kbd{Ctrl+Shift+S} \\
  \menu{Recover Last Session...}\index{recovery!Recover Last Session} & Restore the most recent autosave snapshot. & §\ref{sec:autosave} \\
  \menu{Generate Project Report...}\index{project report} & Create a PDF summary of the dataset and analysis. & §\ref{sec:project-report} \\
  \menu{Export As > Export Figure}\index{export!figure} & Save the plot as an image or vector figure. & \kbd{Ctrl+E}, §\ref{sec:export-figure} \\
  \menu{Export As > Export to FIT}\index{export!FITS} & Save the processed spectrum as a FITS file. & \kbd{Ctrl+F}, §\ref{sec:export-fits} \\
  \menu{Export As > Export Provenance Report...}\index{provenance report} & Write Markdown and JSON provenance summaries. & §\ref{sec:provenance} \\
  \menu{Export As > Export Analysis Log...}\index{analysis log} & Write CSV and text summaries of the fit and shock results. & §\ref{sec:provenance} \\
  \bottomrule
\end{xltabular}

\begin{note}
  The disabled \menu{File > Save As} entry in the main window is reserved; projects are
  saved with \menu{Save Project} and \menu{Save Project As...}.
\end{note}

\begin{xltabular}{\linewidth}{@{}P{0.36\linewidth} L P{0.17\linewidth}@{}}
  \caption{Edit menu.}\label{tab:menu-edit}\\
  \toprule \thead{Command} & \thead{Function} & \thead{Shortcut, see} \\ \midrule\endfirsthead
  \toprule \thead{Command} & \thead{Function} & \thead{Shortcut, see} \\ \midrule\endhead
  \menu{Undo} & Undo the last processing or timeline change. & \kbd{Ctrl+Z} \\
  \menu{Redo} & Redo the last undone change. & \kbd{Ctrl+Shift+Z} \\
  \menu{Add Previous Observation} & Prepend the preceding 15-minute observation. & \kbd{Ctrl+Shift+Left}, §\ref{sec:timeline} \\
  \menu{Add Next Observation} & Append the following observation. & \kbd{Ctrl+Shift+Right}, §\ref{sec:timeline} \\
  \menu{Prefetch Adjacent Observations}\index{timeline!prefetch} & When ticked (default), download neighboring observations in the background so the next timeline step is instant. & §\ref{sec:timeline} \\
  \menu{Reset to Raw} & Remove all processing and return to the loaded data. & \\
  \menu{Reset All} & Clear processing, selections, analysis overlays and the loaded data. & \\
  \bottomrule
\end{xltabular}

\begin{xltabular}{\linewidth}{@{}P{0.36\linewidth} L P{0.17\linewidth}@{}}
  \caption{Download and Solar Events menus.}\label{tab:menu-solar-events}\\
  \toprule \thead{Command} & \thead{Function} & \thead{See} \\ \midrule\endfirsthead
  \toprule \thead{Command} & \thead{Function} & \thead{See} \\ \midrule\endhead
  \menu{Download > Launch FITS Downloader} & Open the e-CALLISTO FITS Downloader. & Ch.~\ref{ch:downloader} \\
  \menu{Solar Events > CMEs > SOHO/LASCO CME Catalog}\index{SOHO/LASCO} & Daily CME lists, parameters and movies. & §\ref{sec:lasco-catalog} \\
  \menu{Solar Events > Flares > GOES X-Ray Flux} & GOES XRS viewer and flare parameters. & §\ref{sec:goes-xrs} \\
  \menu{Solar Events > Energetic Particles > GOES SEP Proton Flux} & GOES SGPS proton flux viewer. & §\ref{sec:goes-sep} \\
  \menu{Solar Events > Geomagnetic > Kyoto Dst Index} & Dst index viewer. & §\ref{sec:dst} \\
  \menu{Solar Events > Geomagnetic > GFZ Kp Index} & Kp index viewer. & §\ref{sec:kp} \\
  \menu{Solar Events > Archives > SunPy Multi-Mission Explorer} & Multi-mission archive search and plotting. & Ch.~\ref{ch:sunpy} \\
  \menu{Solar Events > Radio Bursts > e-CALLISTO} & The e-CALLISTO FITS Downloader. & Ch.~\ref{ch:downloader} \\
  \menu{Solar Events > Radio Bursts > Learmonth} & Learmonth spectrograph browser. & §\ref{sec:learmonth} \\
  \menu{Solar Events > Radio Bursts > SWAVES} & STEREO/SWAVES loader. & §\ref{sec:swaves} \\
  \menu{Solar Events > Sync Current Time Window}\index{time synchronization} & Send the time range of the loaded spectrum to open context windows. & §\ref{sec:sync-time} \\
  \menu{Solar Events > GOES Overlay > Short(XRS-A)}\index{GOES!overlay}, \menu{Long(XRS-B)} & Draw GOES X-ray curves on the spectrum. & §\ref{sec:goes-overlay} \\
  \menu{Solar Events > SWAVES Panel}\index{STEREO/SWAVES!panel} & Show or hide the loaded SWAVES panel. & §\ref{sec:swaves} \\
  \bottomrule
\end{xltabular}

\begin{xltabular}{\linewidth}{@{}P{0.36\linewidth} L P{0.17\linewidth}@{}}
  \caption{View menu.}\label{tab:menu-view}\\
  \toprule \thead{Command} & \thead{Function} & \thead{See} \\ \midrule\endfirsthead
  \toprule \thead{Command} & \thead{Function} & \thead{See} \\ \midrule\endhead
  \menu{View FITS Header} & Show the primary header of the loaded file. & §\ref{sec:fits-header} \\
  \menu{Set Display Range...}\index{display range!dialog} & Enter exact time and frequency limits. & §\ref{sec:display-range} \\
  \menu{Save / Apply / Delete Display Range Preset...}\index{display range!presets} & Store and reuse named display ranges. & §\ref{sec:display-range} \\
  \menu{Export View Config...}, \menu{Import View Config...} & Share display settings as \file{.efaview.json}\index{efaview.json@.efaview.json file}. & §\ref{sec:view-config} \\
  \menu{Multi-Station Comparison...} & Open the comparison workspace. & Ch.~\ref{ch:comparison} \\
  \menu{Theme > System / Light / Dark} & Choose the color theme. & §\ref{sec:themes} \\
  \menu{Mode > Classic / Modern} & Choose the Matplotlib or hardware-accelerated renderer. & §\ref{sec:renderers} \\
  \bottomrule
\end{xltabular}

\begin{xltabular}{\linewidth}{@{}P{0.36\linewidth} L P{0.17\linewidth}@{}}
  \caption{Analysis menu.}\label{tab:menu-analysis}\\
  \toprule \thead{Command} & \thead{Function} & \thead{See} \\ \midrule\endfirsthead
  \toprule \thead{Command} & \thead{Function} & \thead{See} \\ \midrule\endhead
  \menu{Solar Image Analysis} & Open the multi-mission imaging workspace. & Ch.~\ref{ch:solar-window} \\
  \menu{GCS CME Fitting…} & Open the GCS window for the visible time range. & Ch.~\ref{ch:gcs-window} \\
  \menu{Maximum Intensities > Open Maximum Intensities} & Extract the peak frequency of each time channel. & §\ref{sec:max-intensities} \\
  \menu{Type II Band-splitting > Open Type II Band-splitting} & Open the band-splitting analyzer. & Ch.~\ref{ch:band-splitting} \\
  \menu{Plot Light Curves > Enter frequency} & Plot a light curve at a typed frequency. & §\ref{sec:light-curves} \\
  \menu{Plot Light Curves > Click on a frequency}\index{light curve!click mode} & Toggle click mode for light curves. & §\ref{sec:light-curves} \\
  \menu{Plot Light Curves > Settings...} & Light-curve appearance. & §\ref{sec:light-curves} \\
  \menu{Plot Light Curves > Clear light curve(s)} & Remove all light curves. & §\ref{sec:light-curves} \\
  \menu{Ruler Measurement} & Measure duration, frequency change and slope between two points. & §\ref{sec:ruler} \\
  \menu{Clear Ruler Measurement} & Remove the ruler. & §\ref{sec:ruler} \\
  \bottomrule
\end{xltabular}

\begin{xltabular}{\linewidth}{@{}P{0.36\linewidth} L P{0.17\linewidth}@{}}
  \caption{Processing, About and Help menus.}\label{tab:menu-processing}\\
  \toprule \thead{Command} & \thead{Function} & \thead{Shortcut, see} \\ \midrule\endfirsthead
  \toprule \thead{Command} & \thead{Function} & \thead{Shortcut, see} \\ \midrule\endhead
  \menu{Processing > Hardware Acceleration > Enable} & Toggle the accelerated live canvas. & §\ref{sec:renderers} \\
  \menu{Processing > RFI Cleaning > Open RFI Panel / Apply RFI / Reset RFI} & RFI cleaning pipeline. & §\ref{sec:rfi} \\
  \menu{Processing > Annotations > …} & Polygon, line and text annotations. & §\ref{sec:annotations} \\
  \menu{Processing > Presets > …} & Processing presets and the default preset. & §\ref{sec:presets} \\
  \menu{Processing > Maximum Intensity > Auto-Clean Isolated Burst Outliers}\index{outlier removal!automatic} & Remove outliers automatically when maxima are taken from an isolated burst. & §\ref{sec:outliers} \\
  \menu{Processing > Analysis Session > Open Restored Analysis}\index{restored analysis session} & Reopen the analysis windows of a restored session. & §\ref{sec:restored-analysis} \\
  \menu{Processing > Batch Processing > Open Batch Processor} & Export many files with consistent settings. & §\ref{sec:batch} \\
  \menu{About > Check for Updates...}\index{updates!check} & Look for a newer release. & §\ref{sec:updates} \\
  \menu{About > Report a Bug...} & Prepare a bug report with diagnostics. & §\ref{sec:bug-report} \\
  \menu{About > About} & Version and author information. & §\ref{sec:about} \\
  \menu{Help > User Guide} & Open the built-in user guide. & \kbd{F1}, §\ref{sec:in-app-guide} \\
  \bottomrule
\end{xltabular}

\section{Status bar}
\label{sec:status-bar}
\index{status bar}\index{cursor readout}

The left part of the status bar shows messages from the current operation, for example
the result of a drift-rate estimate or a reminder of how to finish a tool. The right part
shows:
\begin{itemize}
  \item the update status, for example \gui{Updates: up to date}; and
  \item the live cursor readout \gui{t = … | f = … MHz | I = …}, giving the time,
    frequency and intensity (in the selected units) under the mouse pointer.
\end{itemize}
The readout allows quick quantitative inspection of the spectrum without clicking.

\section{Themes}
\label{sec:themes}
\index{theme}

\menu{View > Theme}\index{theme} selects \gui{System} (follow the operating system), \gui{Light} or
\gui{Dark}. The theme applies to all windows of the application and to on-screen plots.
Exported figures and PDF reports are always drawn on a white page, whichever theme is in
use.

\section{Drag and drop}
\index{drag and drop}

FITS files and \file{.efaproj} projects can be dropped anywhere on the main window,
including the plot. A single FITS file opens on its own; several are combined exactly as
with \menu{File > Open}; a project file opens the project. Other file types are refused.
If the current work has unsaved changes, you are asked whether to save them first.
